\scanpage{4}\mspage{3}
\[
k(C)\handwritten{\in}H^{\handwritten{3}}(\handwritten{\pi}_{0}(C),\handwritten{\unsure{K|\pi}}_{1}(C))\ .
\]
En précisant ces réflexions, on trouve que la classification, à $\otimes$-équivalence près, des (Gr)-catégories, se fait précisément en termes d\deleted{\rlap{x}\unsure{e}}'un groupe $\handwritten{\pi}_{0}$, d'un groupe commutatif $\handwritten{\pi}_{1}$ sur lequel $\handwritten{\pi}_{0}$ opère, et d'une classe de cohomologie (qui peut être prise arbitraire) dans $H^{\handwritten{3}}(\handwritten{\pi}_{0},\handwritten{\pi}_{1})$. La loi de \deleted{\rlap{xxxxx}comp} groupe du $H^{\handwritten{3}}$ admet d'ailleurs une interprétation ``géométrique'' à la Baer, en termes d'opérations sur les \handwritten{(}Gr)-catégories.

\uline{Cas particuliers}\quad Dans l'exemple 1, on \inserted{doit} retrouver bien sûr le premier invariant de Postnikoff $\handwritten{\in}H^{\handwritten{3}}(\handwritten{\pi}_{1}(X),\handwritten{\pi}_{2}(X))$, avec une interprétation ``géo\deleted{(}métrique'' nouvelle. Dans l'exemple 2, on trouve $\handwritten{\pi}_{0}=H^{1}(X,F)$, $\handwritten{\pi}_{1}=H^{0}(X,F)$\handwritten{,} et l'élément du $H^{\handwritten{3}}$ est nul (du fait qu'on a une ``catégorie de Picard stricte'', cf\inserted{\handwritten{ Deligne}} \handwritten{[\, ]}). Dans \deleted{\rlap{xxx}let} le dernier exemple, on a $\handwritten{\pi}_{0}={}$ groupe des classes d'isomorphisme d'autoéquivalences de $A$, $\handwritten{\pi}_{1}={}$ groupe des automorphismes du foncteur identique de $A$, et l'invariant $k(C)$ semble nouveau (même dans le cas où $C$ est la catégorie des biModules inversibles sous un Anneau $\underline{O}_{X}$, où $\handwritten{\pi}_{0}$ peut donc s'interpréter comme un ``groupoide de Brandt'' et $\handwritten{\pi}_{1}$ comme le groupe des ``unités centrales'' de $\underline{O}_{X}$). Signalons cependant que lorsque $A$ est la catégorie des torseurs sous un groupe (ordinaire) $F$, on trouve \deleted{\rlap{xxxxxxxxxxxxxxxxxxx}un élement de $H^{\handwritten{3}}($} $\handwritten{\pi}_{0}={}$ groupe des automorphismes extérieurs de $F$, $\handwritten{\pi}_{1}={}$ centre de $F$, et l'élément $k(C)$ du $H^{\handwritten{3}}$ n'est autre que le classique invariant de Mac-Lane. Il resterait à \deleted{\rlap{xxxxxxxx}établir}\inserted{développer} un lien \inserted{géométrique} \uline{direct} entre la théorie de\deleted{\unsure{s|x}} Mac-Lane des extensions de groupes non commutatifs\deleted{\unsure{,|s}} et des liens (``kernels''\deleted{\unsure{)|,}} dans sa terminologie) auxquels ils sont associés, et la théorie \deleted{\rlap{xxxxxxxxxxxxxxxx}développées ici,} esquissée ici des (Gr)-catég\deleted{\unsure{rr|or}}ories et de leurs splittages.

2. \uline{Catégories de Picard}.

On appelle ``catégorie de Picard'' une $\otimes$-catégorie ACU (associative-commutative-unitaire , cf. \handwritten{[\, ]}) qui est une (\deleted{E}Gr)-catégorie pour sa \deleted{\rlap{xxxx}(Gr)} $\otimes$-structure AU. On dit que c'est une catégorie de Picard \uline{stricte} si pour tout $X\handwritten{\in}\mathrm{Ob}\ C$, la symétrie de $X\otimes X$ est l'identité. Ces \deleted{\rlap{xxxxx}catég} dernières catégories sont étudiées systématiquement (dans le contexte plus général des ``champs de Picards'') par Deligne \handwritten{[\, ]}.

Soit $C$ une catégorie de Picard. En tant que (Gr)-catégorie, $C$ possède des invariants $\handwritten{\pi}_{0}$ et $\handwritten{\pi}_{1}$, mais dans le cas actuel $\handwritten{\pi}_{0}$ est aussi commutatif et agit trivialement sur $\handwritten{\pi}_{\handwritten{\unsure{0|1}}}$. On fera attention ici qu'une équivalence de (Gr)-catégories entre deux catégories de Picard $C,C'$ n'est pas nécessairement compatible avec les contraintes de commutativité, donc n'est pas nécessairement une équivalence de catégories de
% archive sheet number: 3
